% 星族（综述）
% license CCBYSA3
% type Wiki

本文根据 CC-BY-SA 协议转载翻译自维基百科\href{https://en.wikipedia.org/wiki/Stellar_population}{相关文章}。

\begin{figure}[ht]
\centering
\includegraphics[width=6cm]{./figures/1ba3ca4c285278ba.png}
\caption{艺术家构想的银河系螺旋结构示意图，展示了巴德提出的恒星总体分类。螺旋臂中的蓝色区域由较年轻的Ⅰ族恒星组成，而中央核球中的黄色恒星则是较古老的Ⅱ族恒星。实际上，许多Ⅰ族恒星也与较老的Ⅱ族恒星混合分布在一起。} \label{fig_Stpo_1}
\end{figure}
1944 年，瓦尔特·巴德（Walter Baade）将银河系中的恒星分组划分为不同的恒星族群。在巴德文章的摘要中，他指出这种分类思想最早由扬·奥尔特（Jan Oort）于 1926 年提出。$^{[1]}$

巴德观察到，偏蓝色的恒星与螺旋臂有着很强的关联，而黄色恒星则主要分布在银河系中心核球附近以及球状星团之中。$^{[2]}$ 因此，他提出了两个主要划分：Ⅰ族恒星（Population I）和Ⅱ族恒星（Population II）；1978 年又增加了一个较新的、假设性的划分，即Ⅲ族恒星（Population III）。

在不同恒星族群之间，人们发现其各自观测到的恒星光谱存在显著差异。后来研究表明，这些差异极其重要，并且可能与恒星形成、观测到的运动学特性$^{[3]}$、恒星年龄，甚至螺旋星系和椭圆星系的演化有关。这三种简单的恒星族群在化学组成（即金属丰度）方面对恒星进行了有效划分。$^{[4][5][3]}$ 在天体物理学术语中，“金属”指的是所有比氦更重的元素，其中也包括诸如氧这样的化学非金属元素。$^{[6]}$

按定义而言，每一个恒星族群都表现出这样一种趋势：金属含量越低，恒星年龄越高。因此，宇宙中最早形成的恒星（金属含量极低）被归为Ⅲ族恒星，古老恒星（低金属丰度）被归为Ⅱ族恒星，而较新的恒星（高金属丰度）则被归为Ⅰ族恒星。$^{[7]}$ 太阳被认为是Ⅰ族恒星，是一颗较新的恒星，其金属丰度相对较高，约为 1.4\%。
\subsection{恒星演化}
对恒星光谱的观测表明，比太阳更古老的恒星所含的重元素数量少于太阳。$^{[3]}$ 这立即表明，金属丰度是在恒星世代更替过程中通过恒星核合成这一过程逐步演化而来的。
\subsubsection{第一代恒星的形成}
根据当前的宇宙学模型，在大爆炸中形成的全部物质主要由氢（75\%）和氦（25\%）组成，只有极其微小的一部分由锂、铍等其他轻元素构成。$^{[8]}$ 当宇宙冷却到足够低的温度时，第一代恒星作为Ⅲ族恒星（Population III）诞生，它们不含任何后期污染的重元素。人们推测，这一特性影响了它们的内部结构，使其恒星质量可达到太阳的数百倍。反过来，这些大质量恒星的演化也极为迅速，其核合成过程产生了最初的 26 种元素（即元素周期表中直到铁为止的元素）。$^{[9]}$

许多理论恒星模型表明，大多数高质量的Ⅲ族恒星会迅速耗尽燃料，并很可能以极其高能的对不稳定超新星（pair-instability supernovae）形式爆炸。这些爆炸会将其物质彻底散播开来，把金属抛射进星际介质（ISM），并被后续世代的恒星所吸收。它们的毁灭意味着，银河系中不应还能观测到高质量的Ⅲ族恒星。$^{[10]}$ 然而，在高红移星系中，可能仍能看到一些Ⅲ族恒星，其光起源于宇宙更早期的历史阶段。$^{[11]}$ 科学家已经发现了一颗极端贫金属恒星的证据，这颗恒星略小于太阳，位于银河系螺旋臂中的一个双星系统内。这一发现为观测更古老的恒星打开了可能性。$^{[12]}$

那些质量过大、无法产生对不稳定超新星的恒星，很可能会通过一种称为光致解体（photodisintegration）的过程坍缩成黑洞。在这一过程中，部分物质可能以相对论喷流的形式逃逸，这同样可能将最初的金属散布到宇宙之中。$^{[13][14][a]}$
\subsubsection{已观测恒星的形成}
迄今为止观测到的最古老恒星$^{[10]}$被称为Ⅱ族恒星（Population II），它们具有极低的金属丰度。$^{[16][7]}$ 随着后续世代恒星的诞生，新形成的恒星逐渐变得更加富含金属，因为它们形成所依赖的气体云吸收了由Ⅲ族恒星等前一代恒星制造并释放出的富金属尘埃。

当这些Ⅱ族恒星死亡时，它们通过行星状星云和超新星爆发将富含金属的物质回馈给星际介质，进一步丰富了星云，而新的恒星正是从这些星云中形成的。因此，最年轻的恒星，包括太阳在内，具有最高的金属含量，被称为Ⅰ族恒星（Population I）。
\subsection{瓦尔特·巴德的化学分类}
\subsubsection{Ⅰ族恒星}
\begin{figure}[ht]
\centering
\includegraphics[width=6cm]{./figures/401c0b6650ed6293.png}
\caption{} \label{fig_Stpo_2}
\end{figure}
Ⅰ族恒星（Population I stars）是三种恒星族群中最年轻、金属丰度最高的一类，更常见于银河系的旋臂之中。太阳被认为是一颗中等的Ⅰ族恒星，而类似太阳的 μ Arae 则富含更多的金属。$^{[17]}$（“富金属”一词用于描述金属丰度显著高于太阳、且这种差异不能用测量误差来解释的恒星。）

Ⅰ族恒星通常围绕银河中心做规则的椭圆轨道运动，其相对速度较低。早期曾假设，Ⅰ族恒星较高的金属丰度使其比另外两类恒星更有可能拥有行星系统，因为行星，尤其是类地行星，被认为是通过金属的吸积形成的。$^{[18]}$ 然而，对开普勒空间望远镜数据的观测发现，在具有不同金属丰度的恒星周围都存在小型行星，而只有较大的、可能的气态巨行星才更集中地分布在金属丰度相对较高的恒星周围——这一结果对气态巨行星形成理论具有重要影响。$^{[19]}$ 在中等Ⅰ族恒星与Ⅱ族恒星之间，还存在一个中间的盘族恒星群体。
\subsubsection{Ⅱ族恒星}
\begin{figure}[ht]
\centering
\includegraphics[width=6cm]{./figures/70684b47769d036f.png}
\caption{} \label{fig_Stpo_3}
\end{figure}
\begin{figure}[ht]
\centering
\includegraphics[width=6cm]{./figures/ef10f92a9e0305c0.png}
\caption{} \label{fig_Stpo_4}
\end{figure}